\documentclass[10pt,a4paper]{amsart}
\usepackage[margin=19mm]{geometry}
\usepackage{amsmath,amssymb,mathtools}
\usepackage[scr=rsfs,cal=euler]{mathalfa}
\usepackage{xfrac}
\usepackage[unicode,hidelinks,bookmarks=false]{hyperref}
\newcommand{\ie}{i.e.\ }
\newcommand{\cf}{cf.\ }
\newcommand{\resp}{respectively\ }
\newcommand{\Subsection}{\subsection}
\providecommand{\nobreakdash}{\nobreak\hbox{-}}
\setlength{\emergencystretch}{2em}
\multlinegap0pt
\usepackage{bm}
\usepackage{bbm}
\usepackage{dsfont}
\usepackage{xspace}
\usepackage{cancel}
\usepackage{mathtools}
\usepackage{amsthm}
\makeatletter
\@ifundefined{n}{\newtheorem{n}{N}}{}
\@ifundefined{lemma}{\newtheorem{lemma}[n]{Lemma}}{}
\makeatother
\newcommand\bbS{\mathbb S}
\newcommand\calC{\mathcal C}
\newcommand\calS{\mathcal S}
\newcommand\domain\Omega
\renewcommand\setminus\smallsetminus
\renewcommand\subset\subseteq
\title[Planar percolation and the loop O(n) model]{Planar percolation and the loop O(n) model}
\author{Alexander Glazman, Matan Harel, Nathan Zelesko}
\date{Source: arXiv:2508.20917v3 (CC BY 4.0)}
\begin{document}
\maketitle

\noindent\emph{Source and licence.} Planar percolation and the loop O(n) model. Version arXiv:2508.20917v3 (CC BY 4.0), distributed under the Creative Commons Attribution 4.0 International licence, as recorded in the arXiv licence metadata for this exact version and shown on the abstract page https://arxiv.org/abs/2508.20917v3. Full rights notice: licenses/arxiv-2508.20917-CC-BY-4.0.txt.\par\smallskip

\noindent\textit{Editorial scope.} Counted unit: the Lemma whose source label is \texttt{lem:jordan-domains} in the article (source0.tex lines 602--609, source label \texttt{lem:jordan-domains}), delivered together with the article's own complete original proof of it (source0.tex lines 612--632, one \texttt{proof} environment of 21 lines). No mathematics is written, repaired or re-derived: statement and proof are the article's own text line for line. Required context retained uncounted: none. Every definition, display and label of those lines is retained unchanged. Omitted from this package and not redistributed: 1--601, 610--611, 633--1029. Nothing inside a retained excerpt is omitted or altered: every retained body is delivered exactly as the article's lines are written, so no omission receipt of any kind accompanies this package. Citations: the retained excerpts carry no citation command at all, so no bibliography entry of the article is transcribed and no reference is redistributed. Reference adaptation: none was needed and none was made; every reference inside the delivered unit resolves inside it, so no unresolved reference or citation is produced. Printed slips kept literal: no printed word, symbol, hypothesis or number of the counted statement or of its proof was altered. Layout and class: the article's own class is \texttt{amsart} with packages including amssymb, mathtools, amsthm, geometry, fontenc, lmodern, inputenc, microtype, enumitem, caption, dsfont, subcaption, bbm, tikz; the wrapper is the lane's pinned \texttt{amsart} editorial wrapper at 10pt on A4, which restates the theorem-like environments the retained text needs and adds the article's own macro definitions that the retained text needs (\texttt{\textbackslash bbS}, \texttt{\textbackslash calC}, \texttt{\textbackslash calS}, \texttt{\textbackslash domain}, \texttt{\textbackslash setminus}, \texttt{\textbackslash subset}), with extra wrapper packages (bm, bbm, dsfont, xspace, cancel, mathtools). The delivered numbering is the unit's own. Assets: no retained span contains a figure environment, a graphics-inclusion command, a PDF-page inclusion, a table or a TikZ picture; no image file is copied into this package, the package declares no asset dependency and no image file is redistributed with it. Licence: version arXiv:2508.20917v3 of this article is distributed under the Creative Commons Attribution 4.0 International licence, as recorded in the arXiv licence metadata for this exact version and shown on the abstract page https://arxiv.org/abs/2508.20917v3. A full rights notice is delivered at licenses/arxiv-2508.20917-CC-BY-4.0.txt. Characters: every retained excerpt is pure ASCII, so the delivered engine, XeLaTeX with its default Latin Modern text font, reports no missing character.
\begin{lemma}\label{lem:jordan-domains}
	The sequence of pairs of sets $\{\domain_n, \calS_n\}$ satisfies the following properties:
	\begin{enumerate}
	\item $\domain_n \subseteq \domain_{n+1}$, $\calS_n\cap \calS_{n+1} = \emptyset$, and $G \subseteq \bigcup_{n} \domain_n$,
 	\item $\domain_n$ is a closed, simply connected subset of~$\bbS^2$,  its boundary $J_n$ contains $\calS_n$ and can be continuously parametrized in such a way that no vertex of~$\calS_n$ is covered twice,
	\item for any vertex~$v \in \Omega_n$, any infinite connected component of $\{\sigma=1\}$ whose closure contains both $v$ and $a$ must contain a vertex in $\calS_n$.
\end{enumerate}
\end{lemma}  

\begin{proof}
	For Item~$1$, it is clear that $\Lambda_n(\rho) \subseteq \Lambda_{n+1}(\rho)$.
	This implies $Q_n \supseteq Q_{n+1}$ and~$\domain_{n}\subseteq \domain_{n+1}$ as requested. 
	Also, $\calS_n\subseteq \partial \Lambda_n(\rho)$ and~$\partial \Lambda_n(\rho) \cap \partial \Lambda_{n+1}(\rho) = \emptyset$, whence~$\calS_n\cap \calS_{n+1} = \emptyset$.
	Since $G$ is connected, the combinatorial balls~$\Lambda_n(\rho)\subseteq \domain_n$ exhaust $G$, whence~$G \subseteq \bigcup_n \domain_n$.
	
	For Item~$2$, note that $Q_n$ is open and path connected, whence $\domain_n$ is closed and simply connected. 
	Their common boundary~$J_n$ is the union of all (finitely many) edges that border~$Q_n$.
	By picking a starting point on $J_n$ and traversing this boundary in a clockwise manner, one obtains acontinuous parametrization $\gamma_n: \bbS^1 \to J_n$.  Given distinct~$x,y\in \bbS^1$, define~$(x,y)$ as the clockwise arc excluding the endpoints. If $\gamma_n(x)=\gamma_n(y)=v\in V$, then both~$\gamma_n((x,y))\cap V$ and~$\gamma_n((y,x))\cap V$ are nonempty (since~$G$ is simple).
	Now, define~$\Lambda_n(\rho)^*$ as the dual graph of~$\Lambda_n(\rho)$ with a vertex~$q_n^*$ corresponding to~$Q_n$.
	If the first and last edges in~$\gamma_n((x,y))$ coincide, then define~$\calC$ as the self-loop in~$\Lambda_n(\rho)^*$  at~$q_n^*$ crossing this edge.  
	Otherwise, denote these edges by~$e_1,e_2$, their duals by~$e_1^*,e_2^*$ and construct~$\calC$ as follows:  start at~$q_n^*$, trace a path along $e_1^*$ and $e_1$ to $v$, then along $e_2$ and $e_2^*$ back to~$q_n^*$.
	In either case, $\calC$ is a Jordan curve that separates~$\gamma_n((x,y))$ from~$\gamma_n((y,x))$ in~$\Lambda_n(\rho) \setminus \{v\}$.
	Thus, $v$ is a cut-vertex of~$\Lambda_n(\rho)$.
	In particular, $v\not\in \calS_n\subset \partial\Lambda_n(\rho)$ since no vertex in~$\partial\Lambda_n(\rho)$ can be a cut-vertex.
	
	For Item~$3$, we can find~$\varepsilon>0$ such that ball of radius~$\varepsilon$ around~$a$ does not intersect~$\domain_n$.
	By our definition of connectivity, there exists a vertex~$u$ in this ball that is connected to~$v$ in~$\{\sigma=1\}$.
	Follow this path from~$v$ to~$u$ and let~$w$ be the last vertex on the path that is contained in~$\domain_n$.
	Then, $w\in J_n$, and hence~$w\in \calS_n$.
\end{proof}

\end{document}
